\documentclass[11pt,letterpaper]{article}
\usepackage[margin=1in,headheight=14pt,headsep=18pt]{geometry}
\usepackage{fontspec}
\usepackage[UTF8,scheme=plain,fontset=fandol]{ctex}
\setmainfont{texgyretermes}[Extension=.otf,UprightFont=*-regular,BoldFont=*-bold,ItalicFont=*-italic,BoldItalicFont=*-bolditalic]
\usepackage{setspace}
\usepackage{graphicx}
\usepackage{longtable,booktabs,array}
\usepackage{calc}
\usepackage{amsmath,amssymb}
\usepackage{fancyhdr}
\usepackage{titlesec}
\usepackage{needspace}
\usepackage{flafter}
\usepackage{placeins}
\usepackage{xurl}
\usepackage[unicode,hidelinks,breaklinks]{hyperref}
\urlstyle{same}
\hypersetup{pdfauthor={},pdftitle={补充方法与结果},pdfsubject={Journal-style computational reanalysis draft}}
\pagestyle{fancy}
\fancyhf{}
\fancyfoot[R]{\thepage}
\renewcommand{\headrulewidth}{0pt}
\setlength{\parindent}{2em}
\setlength{\parskip}{3pt plus 1pt}
\setlength{\emergencystretch}{3em}
\setcounter{secnumdepth}{0}
\titleformat{\section}{\normalfont\large\bfseries}{}{0pt}{}
\titleformat{\subsection}{\normalfont\normalsize\bfseries}{}{0pt}{}
\titleformat{\subsubsection}{\normalfont\normalsize\bfseries\itshape}{}{0pt}{}
\titlespacing*{\section}{0pt}{13pt}{5pt}
\titlespacing*{\subsection}{0pt}{10pt}{3pt}
\titlespacing*{\subsubsection}{0pt}{8pt}{3pt}
\providecommand{\tightlist}{\setlength{\itemsep}{0pt}\setlength{\parskip}{0pt}}
\providecommand{\pandocbounded}[1]{#1}
\newcommand{\APAref}{\par\noindent\hangindent=0.5in\hangafter=1\relax}
\setlength{\LTpre}{4pt}
\setlength{\LTpost}{4pt}
\setlength{\tabcolsep}{4pt}
\setlength{\textfloatsep}{12pt plus 2pt minus 2pt}
\setlength{\intextsep}{10pt plus 2pt minus 2pt}
\renewcommand{\topfraction}{0.85}
\renewcommand{\bottomfraction}{0.75}
\renewcommand{\textfraction}{0.12}
\renewcommand{\floatpagefraction}{0.72}
\setstretch{1.30}
\widowpenalty=10000
\clubpenalty=10000
\begin{document}
\begin{center}\fontsize{16}{20}\selectfont\bfseries 补充方法与结果\end{center}
\section{分析来源与缺失输入}\label{ux5206ux6790ux6765ux6e90ux4e0eux7f3aux5931ux8f93ux5165}

本补充材料配套于参与者层面的微博数据再分析。已知所提供的原研究已经发表，但本次未检索或使用其论文正文、摘要和结果。数据采集细节均据随附方法说明报告。原始输入及其校验值、逐行纳入决定、全部参数候选、警告和失败运行、拟合模型、预测值、重抽样代码、图表及源文档均保存在研究包中。

仅凭小数值无法确定问卷变换。未开展原始 PANAS／SWLS 单位换算、条目信度、奇偶文本信度、词云、人口学效度或时间截止点重建。元数据表明，1,027 名主分析参与者的文本覆盖范围延伸至问卷完成之后，因此这些结果不能确立前瞻性预测。34 名参与者的 ind 值超过一，另有 13 个特征代码缺少映射。未根据分布开展事后清洗。

\section{具体模型设置与搜索记录}\label{ux5177ux4f53ux6a21ux578bux8bbeux7f6eux4e0eux641cux7d22ux8bb0ux5f55}

仅在开发集上使用的参数网格包含 42 种配置和 210 次复合折内拟合。独立结局分别选择参数，共享模型选择一组设置。所有拟合变换都仅使用相应折的训练行。候选检查点保留了分数、警告、状态和折外（OOF）数组，并另存明确的 UID 行序文件。最初采用向量缩放的预检实现，以及修正后的标量缩放最终实现均被保留。修正由留出评估前的合成数值一致性检查触发，参数网格和参与者划分没有变化。

\begingroup\singlespacing\fontsize{10}{12}\selectfont\setlength{\parindent}{0pt}

\begin{longtable}[]{@{}
  >{\raggedright\arraybackslash}p{(\linewidth - 4\tabcolsep) * \real{0.1900}}
  >{\raggedright\arraybackslash}p{(\linewidth - 4\tabcolsep) * \real{0.6600}}
  >{\raggedright\arraybackslash}p{(\linewidth - 4\tabcolsep) * \real{0.1500}}@{}}
\toprule\noalign{}
\begin{minipage}[b]{\linewidth}\raggedright
模型类别
\end{minipage} & \begin{minipage}[b]{\linewidth}\raggedright
选定参数
\end{minipage} & \begin{minipage}[b]{\linewidth}\raggedright
CV 归一化 RMSE
\end{minipage} \\
\midrule\noalign{}
\endhead
\bottomrule\noalign{}
\endlastfoot
训练均值 & \{\} & 1.0016 \\
岭回归 & {[}\{"alpha": 1000.0\}, \{"alpha": 10.0\}, \{"alpha": 10.0\}{]} & 0.9969 \\
独立 RF & {[}\{"max\_\allowbreak{}features": 1.0, "min\_\allowbreak{}samples\_\allowbreak{}leaf": 15\}, \{"max\_\allowbreak{}features": 0.5, "min\_\allowbreak{}samples\_\allowbreak{}leaf": 15\}, \{"max\_\allowbreak{}features": 1.0, "min\_\allowbreak{}samples\_\allowbreak{}leaf": 15\}{]} & 0.9965 \\
独立 ET & {[}\{"max\_\allowbreak{}features": 0.5, "min\_\allowbreak{}samples\_\allowbreak{}leaf": 15\}, \{"max\_\allowbreak{}features": 1.0, "min\_\allowbreak{}samples\_\allowbreak{}leaf": 15\}, \{"max\_\allowbreak{}features": 1.0, "min\_\allowbreak{}samples\_\allowbreak{}leaf": 15\}{]} & 0.9968 \\
独立 HGB & {[}\{"l2\_\allowbreak{}regularization": 10.0, "max\_\allowbreak{}leaf\_\allowbreak{}nodes": 15\}, \{"l2\_\allowbreak{}regularization": 1.0, "max\_\allowbreak{}leaf\_\allowbreak{}nodes": 31\}, \{"l2\_\allowbreak{}regularization": 10.0, "max\_\allowbreak{}leaf\_\allowbreak{}nodes": 31\}{]} & 1.0226 \\
共享 RF & \{"max\_\allowbreak{}features": 1.0, "min\_\allowbreak{}samples\_\allowbreak{}leaf": 15\} & 0.9963 \\
共享 ET & \{"max\_\allowbreak{}features": 0.5, "min\_\allowbreak{}samples\_\allowbreak{}leaf": 15\} & 0.9977 \\
多任务弹性网 & \{"alpha": 0.01, "l1\_\allowbreak{}ratio": 0.8\} & 0.9972 \\
\end{longtable}

\nopagebreak

注. 所有森林均使用 256 棵树。RF 对树使用 bootstrap 抽样；ET 默认使用完整训练样本。HGB 设置为 max\_\allowbreak{}iter = 150、learning\_\allowbreak{}rate = .05、min\_\allowbreak{}samples\_\allowbreak{}leaf = 20、early\_\allowbreak{}stopping = False。弹性网设置为 max\_\allowbreak{}iter = 10000、tol = 1e-5。岭回归和独立集成模型的参数依次对应消极情绪、积极情绪和生活满意度。实际安装库的确切默认设置保存在拟合模型及 environment/requirements.lock.txt 中。

\endgroup

\section{开发集重复交叉验证}\label{ux5f00ux53d1ux96c6ux91cdux590dux4ea4ux53c9ux9a8cux8bc1}

使用随机种子 20261023、20261024 和 20261025，进行了三次打乱顺序的五折重复，重新拟合已固定的选定配置。超参数未重新调节。折间变异属于描述性结果，不是无偏性能的置信区间，因为同一开发样本此前已参与参数选择。

\begingroup\singlespacing\fontsize{10}{12}\selectfont\setlength{\parindent}{0pt}

\begin{longtable}[]{@{}
  >{\raggedright\arraybackslash}p{(\linewidth - 6\tabcolsep) * \real{0.2300}}
  >{\raggedright\arraybackslash}p{(\linewidth - 6\tabcolsep) * \real{0.2600}}
  >{\raggedright\arraybackslash}p{(\linewidth - 6\tabcolsep) * \real{0.2800}}
  >{\raggedright\arraybackslash}p{(\linewidth - 6\tabcolsep) * \real{0.2300}}@{}}
\toprule\noalign{}
\begin{minipage}[b]{\linewidth}\raggedright
模型
\end{minipage} & \begin{minipage}[b]{\linewidth}\raggedright
结局
\end{minipage} & \begin{minipage}[b]{\linewidth}\raggedright
平均归一化 RMSE
\end{minipage} & \begin{minipage}[b]{\linewidth}\raggedright
折间 SD
\end{minipage} \\
\midrule\noalign{}
\endhead
\bottomrule\noalign{}
\endlastfoot
共享 RF & 消极情绪 & 0.9993 & 0.0627 \\
共享 RF & 积极情绪 & 0.9995 & 0.0523 \\
共享 RF & 生活满意度 & 0.9923 & 0.0578 \\
独立 RF & 消极情绪 & 1.0034 & 0.0647 \\
独立 RF & 积极情绪 & 1.0003 & 0.0548 \\
独立 RF & 生活满意度 & 0.9899 & 0.0615 \\
训练均值 & 消极情绪 & 1.0005 & 0.0647 \\
训练均值 & 积极情绪 & 1.0011 & 0.0568 \\
训练均值 & 生活满意度 & 1.0007 & 0.0570 \\
\end{longtable}

\endgroup

\section{完整留出测试集表现}\label{ux5b8cux6574ux7559ux51faux6d4bux8bd5ux96c6ux8868ux73b0}

所有模型均使用同一批 104 名参与者，这些参与者未参与模型选择。下表中的 Pearson P 值未经校正；对非主要模型而言，这些结果属于描述性信息。主要选定模型的三个结局 P 值已在正文结果中进行 Holm 校正。常量预测的相关系数无定义。采用 CSV 单位的 RMSE 区间以拟合模型及选定的训练／留出划分为条件。

\Needspace{8\baselineskip}

\par\medskip\noindent{\bfseries 表 S1 各模型在主分析留出测试集上的表现}\par\nopagebreak

\begingroup\singlespacing\fontsize{10}{12}\selectfont\setlength{\parindent}{0pt}

\begin{longtable}[]{@{}
  >{\raggedright\arraybackslash}p{(\linewidth - 8\tabcolsep) * \real{0.1700}}
  >{\raggedright\arraybackslash}p{(\linewidth - 8\tabcolsep) * \real{0.1900}}
  >{\raggedright\arraybackslash}p{(\linewidth - 8\tabcolsep) * \real{0.2500}}
  >{\raggedright\arraybackslash}p{(\linewidth - 8\tabcolsep) * \real{0.1000}}
  >{\raggedright\arraybackslash}p{(\linewidth - 8\tabcolsep) * \real{0.2900}}@{}}
\toprule\noalign{}
\begin{minipage}[b]{\linewidth}\raggedright
模型
\end{minipage} & \begin{minipage}[b]{\linewidth}\raggedright
结局
\end{minipage} & \begin{minipage}[b]{\linewidth}\raggedright
Pearson r［95\% CI］
\end{minipage} & \begin{minipage}[b]{\linewidth}\raggedright
p
\end{minipage} & \begin{minipage}[b]{\linewidth}\raggedright
RMSE［95\% CI］
\end{minipage} \\
\midrule\noalign{}
\endhead
\bottomrule\noalign{}
\endlastfoot
训练均值 & 消极情绪 & 无定义 & --- & 0.169 {[}0.150, 0.190{]} \\
训练均值 & 积极情绪 & 无定义 & --- & 0.170 {[}0.146, 0.192{]} \\
训练均值 & 生活满意度 & 无定义 & --- & 0.205 {[}0.184, 0.225{]} \\
岭回归 & 消极情绪 & 0.274 {[}0.086, 0.443{]} & 0.005 & 0.169 {[}0.148, 0.189{]} \\
岭回归 & 积极情绪 & 0.330 {[}0.147, 0.492{]} & \textless{} .001 & 0.163 {[}0.140, 0.186{]} \\
岭回归 & 生活满意度 & 0.166 {[}\mbox{-0.027}, 0.348{]} & 0.092 & 0.202 {[}0.181, 0.222{]} \\
独立 RF & 消极情绪 & 0.230 {[}0.039, 0.405{]} & 0.019 & 0.165 {[}0.144, 0.185{]} \\
独立 RF & 积极情绪 & 0.325 {[}0.141, 0.487{]} & \textless{} .001 & 0.163 {[}0.141, 0.186{]} \\
独立 RF & 生活满意度 & 0.184 {[}\mbox{-0.008}, 0.364{]} & 0.061 & 0.201 {[}0.180, 0.222{]} \\
独立 ET & 消极情绪 & 0.276 {[}0.088, 0.445{]} & 0.005 & 0.165 {[}0.145, 0.186{]} \\
独立 ET & 积极情绪 & 0.275 {[}0.087, 0.444{]} & 0.005 & 0.165 {[}0.142, 0.187{]} \\
独立 ET & 生活满意度 & 0.225 {[}0.033, 0.400{]} & 0.022 & 0.200 {[}0.179, 0.220{]} \\
独立 HGB & 消极情绪 & 0.075 {[}\mbox{-0.119}, 0.264{]} & 0.447 & 0.172 {[}0.150, 0.194{]} \\
独立 HGB & 积极情绪 & 0.219 {[}0.028, 0.395{]} & 0.025 & 0.167 {[}0.146, 0.187{]} \\
独立 HGB & 生活满意度 & 0.190 {[}\mbox{-0.003}, 0.369{]} & 0.054 & 0.204 {[}0.180, 0.227{]} \\
共享 RF & 消极情绪 & 0.247 {[}0.057, 0.419{]} & 0.012 & 0.166 {[}0.146, 0.186{]} \\
共享 RF & 积极情绪 & 0.284 {[}0.097, 0.452{]} & 0.003 & 0.165 {[}0.143, 0.187{]} \\
共享 RF & 生活满意度 & 0.152 {[}\mbox{-0.042}, 0.335{]} & 0.124 & 0.202 {[}0.181, 0.222{]} \\
共享 ET & 消极情绪 & 0.312 {[}0.127, 0.476{]} & 0.001 & 0.166 {[}0.145, 0.186{]} \\
共享 ET & 积极情绪 & 0.306 {[}0.120, 0.471{]} & 0.002 & 0.166 {[}0.142, 0.188{]} \\
共享 ET & 生活满意度 & 0.227 {[}0.036, 0.402{]} & 0.021 & 0.201 {[}0.180, 0.222{]} \\
多任务弹性网 & 消极情绪 & 0.331 {[}0.148, 0.493{]} & \textless{} .001 & 0.166 {[}0.146, 0.186{]} \\
多任务弹性网 & 积极情绪 & 0.402 {[}0.227, 0.552{]} & \textless{} .001 & 0.165 {[}0.143, 0.187{]} \\
多任务弹性网 & 生活满意度 & 0.168 {[}\mbox{-0.025}, 0.350{]} & 0.087 & 0.202 {[}0.182, 0.223{]} \\
\end{longtable}

\endgroup

\section{固定参数敏感性分析}\label{ux56faux5b9aux53c2ux6570ux654fux611fux6027ux5206ux6790}

最新记录分析保留同样的 1,037 名参与者及其划分，使用数值上最新的问卷时间。单次 ID 分析排除所有重复 ID，因此样本量减小。受限特征分析保留参与者，但去除 ind 和 A、B、C1--C10、D，剩余 106 项预测变量。每项敏感性分析均重新拟合各模型原先选定的超参数，不重新调参。全部模型与情境组合的结果，包括决定系数和 Spearman 相关，保存在 tables/tableS2\_\allowbreak{}sensitivity.csv。下表重点呈现原先选定的模型，并未根据敏感性结果重新选出模型。

\Needspace{8\baselineskip}

\par\medskip\noindent{\bfseries 表 S2 原选定模型在预定数据集上的表现}\par\nopagebreak

\begingroup\singlespacing\fontsize{10}{12}\selectfont\setlength{\parindent}{0pt}

\begin{longtable}[]{@{}
  >{\raggedright\arraybackslash}p{(\linewidth - 10\tabcolsep) * \real{0.1900}}
  >{\raggedright\arraybackslash}p{(\linewidth - 10\tabcolsep) * \real{0.2000}}
  >{\raggedright\arraybackslash}p{(\linewidth - 10\tabcolsep) * \real{0.0800}}
  >{\raggedright\arraybackslash}p{(\linewidth - 10\tabcolsep) * \real{0.0800}}
  >{\raggedright\arraybackslash}p{(\linewidth - 10\tabcolsep) * \real{0.3000}}
  >{\raggedright\arraybackslash}p{(\linewidth - 10\tabcolsep) * \real{0.1500}}@{}}
\toprule\noalign{}
\begin{minipage}[b]{\linewidth}\raggedright
数据集
\end{minipage} & \begin{minipage}[b]{\linewidth}\raggedright
结局
\end{minipage} & \begin{minipage}[b]{\linewidth}\raggedright
测试集 n
\end{minipage} & \begin{minipage}[b]{\linewidth}\raggedright
r
\end{minipage} & \begin{minipage}[b]{\linewidth}\raggedright
RMSE［95\% CI］
\end{minipage} & \begin{minipage}[b]{\linewidth}\raggedright
决定系数
\end{minipage} \\
\midrule\noalign{}
\endhead
\bottomrule\noalign{}
\endlastfoot
主分析 & 消极情绪 & 104 & 0.247 & 0.166 {[}0.146, 0.186{]} & 0.020 \\
主分析 & 积极情绪 & 104 & 0.284 & 0.165 {[}0.143, 0.187{]} & 0.040 \\
主分析 & 生活满意度 & 104 & 0.152 & 0.202 {[}0.181, 0.222{]} & 0.015 \\
最新记录 & 消极情绪 & 104 & 0.259 & 0.174 {[}0.154, 0.194{]} & 0.024 \\
最新记录 & 积极情绪 & 104 & 0.286 & 0.163 {[}0.144, 0.182{]} & 0.048 \\
最新记录 & 生活满意度 & 104 & 0.101 & 0.207 {[}0.185, 0.228{]} & 0.008 \\
单次 ID & 消极情绪 & 85 & 0.303 & 0.171 {[}0.147, 0.195{]} & 0.025 \\
单次 ID & 积极情绪 & 85 & 0.301 & 0.161 {[}0.138, 0.184{]} & 0.055 \\
单次 ID & 生活满意度 & 85 & 0.173 & 0.207 {[}0.185, 0.229{]} & 0.027 \\
受限特征 & 消极情绪 & 104 & 0.301 & 0.165 {[}0.146, 0.185{]} & 0.031 \\
受限特征 & 积极情绪 & 104 & 0.268 & 0.165 {[}0.142, 0.188{]} & 0.035 \\
受限特征 & 生活满意度 & 104 & 0.125 & 0.203 {[}0.182, 0.223{]} & 0.008 \\
\end{longtable}

\endgroup

\section{归因与置换诊断}\label{ux5f52ux56e0ux4e0eux7f6eux6362ux8bcaux65ad}

TreeSHAP 使用在 933 名参与者上拟合的共享 RF。用期望值加上特征归因重建训练预测时，最大误差为 2.22e-16 个 CSV 单位。各目标训练重要性前十的特征取并集后，共得到 14 个不同特征。在 104 名留出参与者中逐一置换这些特征，每个特征洗牌 20 次。平均 RMSE 差值及洗牌 SD 衡量的是条件性的模型敏感性，其不确定性不是总体置信区间，相关类别还可能相互替代。完整的带符号贡献、期望值、输入值和参与者标识保存在 results/primary/interpretation/training\_\allowbreak{}shap\_\allowbreak{}values.npz。

\begingroup\singlespacing\fontsize{10}{12}\selectfont\setlength{\parindent}{0pt}

\begin{longtable}[]{@{}
  >{\raggedright\arraybackslash}p{(\linewidth - 6\tabcolsep) * \real{0.2800}}
  >{\raggedright\arraybackslash}p{(\linewidth - 6\tabcolsep) * \real{0.2800}}
  >{\raggedright\arraybackslash}p{(\linewidth - 6\tabcolsep) * \real{0.2300}}
  >{\raggedright\arraybackslash}p{(\linewidth - 6\tabcolsep) * \real{0.2100}}@{}}
\toprule\noalign{}
\begin{minipage}[b]{\linewidth}\raggedright
结局
\end{minipage} & \begin{minipage}[b]{\linewidth}\raggedright
训练阶段选定特征
\end{minipage} & \begin{minipage}[b]{\linewidth}\raggedright
平均 RMSE 差值
\end{minipage} & \begin{minipage}[b]{\linewidth}\raggedright
洗牌 SD
\end{minipage} \\
\midrule\noalign{}
\endhead
\bottomrule\noalign{}
\endlastfoot
消极情绪 & Sexual & 0.00008 & 0.00052 \\
消极情绪 & C10 & 0.00032 & 0.00050 \\
消极情绪 & C2 & 0.00025 & 0.00044 \\
消极情绪 & D & 0.00007 & 0.00038 \\
消极情绪 & Angery & 0.00011 & 0.00029 \\
积极情绪 & Sexual & 0.00000 & 0.00054 \\
积极情绪 & We & 0.00023 & 0.00041 \\
积极情绪 & Family & \mbox{-0.00008} & 0.00034 \\
积极情绪 & D & 0.00022 & 0.00030 \\
积极情绪 & Angery & 0.00023 & 0.00030 \\
生活满意度 & Family & 0.00103 & 0.00069 \\
生活满意度 & We & \mbox{-0.00093} & 0.00074 \\
生活满意度 & C10 & 0.00006 & 0.00063 \\
生活满意度 & Sexual & \mbox{-0.00017} & 0.00066 \\
生活满意度 & D & 0.00093 & 0.00058 \\
\end{longtable}

\endgroup

\section{独立计算复现}\label{ux72ecux7acbux8ba1ux7b97ux590dux73b0}

最终分析在锁定环境下，由 data.csv 在独立输出目录中重新运行。参与者清单、CV 行序、选定参数、CV 折分数、主分析预测值及已报告的主分析指标，均在 1e-12 的容差内一致。主分析预测值的最大绝对差为 0.0e+00。这是在相同输入和环境下的计算可复现性，而不是独立样本重复研究。JSON 核验报告列出了确切检查项，以及任何有意未重复的非主要情境。固定配置的开发集重复交叉验证保留自原最终运行；验证主分析预测等价性无需再次执行该诊断。

\section{数据与文档核验}\label{ux6570ux636eux4e0eux6587ux6863ux6838ux9a8c}

终稿参考文献包含 30 项学术文献，均发表于 2024 年 8 月 1 日之前。全部保存在 Zotero 集合 S2S9AGIR，具有可读全文附件及记录在案的原文位置。新增 16 篇均通过 Zotero MCP 入库并回读；所有导入的本地文件 SHA-256 均与保留的来源文件一致。集合共有 31 个来源附件，包括此前 Lundberg 文本与预印本这一组附件。Diener 等（2018）以 NCBI 作者稿纯文本 RTF 保存，其中行内引文标记和图形并不完整。Diener 等（1999）的扫描原文配有本地 OCR，用于核对段落。证据库明确标注版本差异与限制。未检索、获取、使用或引用目标研究的已发表原文。期刊规则依据 2026 年 9 月 23 日核查的官方来源。作者控制的声明和负责任的人类审核仍待完成。

\section{主分析后的探索性优化}\label{ux4e3bux5206ux6790ux540eux7684ux63a2ux7d22ux6027ux4f18ux5316}

包含 22 个配置的搜索是在查看原始留出集之后固定的。它仍使用相同的 933 名开发参与者和 104 名留出参与者。五个开发集外层折各自执行内层五折选择，之后在全部开发数据上进行五折搜索，共形成 132 条候选配置与搜索组合记录及 660 次复合折评估。所有评估均完成，未出现失败或捕获到的警告。除明确共享输出的候选外，标量结局估计器分别拟合。岭回归、支持向量回归、弹性网与浅层提升考察不同的平滑或正则化假设；PLS 考察低维协变结构。所有需要拟合的插补器、缩放器及结局变换仅使用对应训练折。固定的 log1p 变换仅应用于非负输入特征。这是探索性扩展，并非新获得的独立确认数据集。

外层折中，原共享森林被选择一次，alpha 为 0.1、L1 比例为 0.5 的弹性网被选择两次，叶节点样本量为 60 的独立森林被选择两次。这些选择差异作为诊断信息保留。表 S3 给出五个外层折的未加权均值及描述性 SD；这些 SD 不是标准误或置信区间，训练集重叠的各折也不视为独立重复。

\Needspace{8\baselineskip}

\par\medskip\noindent{\bfseries 表 S3 嵌套开发集验证}\par\nopagebreak

\begingroup\singlespacing\fontsize{10}{12}\selectfont\setlength{\parindent}{0pt}

\begin{longtable}[]{@{}
  >{\raggedright\arraybackslash}p{(\linewidth - 6\tabcolsep) * \real{0.2800}}
  >{\raggedright\arraybackslash}p{(\linewidth - 6\tabcolsep) * \real{0.2500}}
  >{\raggedright\arraybackslash}p{(\linewidth - 6\tabcolsep) * \real{0.2900}}
  >{\raggedright\arraybackslash}p{(\linewidth - 6\tabcolsep) * \real{0.1800}}@{}}
\toprule\noalign{}
\begin{minipage}[b]{\linewidth}\raggedright
流程
\end{minipage} & \begin{minipage}[b]{\linewidth}\raggedright
结局
\end{minipage} & \begin{minipage}[b]{\linewidth}\raggedright
平均归一化 RMSE
\end{minipage} & \begin{minipage}[b]{\linewidth}\raggedright
折间 SD
\end{minipage} \\
\midrule\noalign{}
\endhead
\bottomrule\noalign{}
\endlastfoot
选定流程 & 消极情绪 & 1.0044 & 0.0481 \\
选定流程 & 积极情绪 & 0.9997 & 0.0822 \\
选定流程 & 生活满意度 & 0.9976 & 0.0572 \\
原共享 RF & 消极情绪 & 1.0027 & 0.0473 \\
原共享 RF & 积极情绪 & 0.9986 & 0.0857 \\
原共享 RF & 生活满意度 & 0.9991 & 0.0578 \\
均值基线 & 消极情绪 & 1.0020 & 0.0463 \\
均值基线 & 积极情绪 & 1.0002 & 0.0836 \\
均值基线 & 生活满意度 & 1.0001 & 0.0605 \\
\end{longtable}

\endgroup

注。开发集外层五折；全部 933 名开发参与者在每种流程下各获得一次外层折预测。折分数以结局 RMSE 除以相应外层训练折的结局 SD。选定流程在每个外层训练集内重新选择候选，另外两个参照使用固定设置。

\Needspace{8\baselineskip}

\par\medskip\noindent{\bfseries 表 S4 完整探索性配置搜索}\par\nopagebreak

\begingroup\singlespacing\fontsize{10}{12}\selectfont\setlength{\parindent}{0pt}

\begin{longtable}[]{@{}
  >{\raggedright\arraybackslash}p{(\linewidth - 4\tabcolsep) * \real{0.2600}}
  >{\raggedright\arraybackslash}p{(\linewidth - 4\tabcolsep) * \real{0.5000}}
  >{\raggedright\arraybackslash}p{(\linewidth - 4\tabcolsep) * \real{0.2400}}@{}}
\toprule\noalign{}
\begin{minipage}[b]{\linewidth}\raggedright
候选
\end{minipage} & \begin{minipage}[b]{\linewidth}\raggedright
参数
\end{minipage} & \begin{minipage}[b]{\linewidth}\raggedright
开发集 CV 分数
\end{minipage} \\
\midrule\noalign{}
\endhead
\bottomrule\noalign{}
\endlastfoot
mean & 固定参照 & 1.0007 \\
original\_\allowbreak{}rf & 固定参照 & 1.0001 \\
ridge\_\allowbreak{}raw\_\allowbreak{}10 & \{"log": false, "alpha": 10\} & 1.0703 \\
ridge\_\allowbreak{}raw\_\allowbreak{}100 & \{"log": false, "alpha": 100\} & 1.0298 \\
ridge\_\allowbreak{}raw\_\allowbreak{}1000 & \{"log": false, "alpha": 1000\} & 1.0013 \\
ridge\_\allowbreak{}log\_\allowbreak{}10 & \{"log": true, "alpha": 10\} & 1.0709 \\
ridge\_\allowbreak{}log\_\allowbreak{}100 & \{"log": true, "alpha": 100\} & 1.0301 \\
ridge\_\allowbreak{}log\_\allowbreak{}1000 & \{"log": true, "alpha": 1000\} & 1.0013 \\
svr\_\allowbreak{}0.1 & \{"C": 0.1\} & 1.0040 \\
svr\_\allowbreak{}1 & \{"C": 1\} & 1.0215 \\
svr\_\allowbreak{}10 & \{"C": 10\} & 1.1019 \\
hgb\_\allowbreak{}10 & \{"l2": 10\} & 1.0139 \\
hgb\_\allowbreak{}100 & \{"l2": 100\} & 1.0013 \\
en\_\allowbreak{}0.01\_\allowbreak{}0.1 & \{"alpha": 0.01, "l1": 0.1\} & 1.0654 \\
en\_\allowbreak{}0.01\_\allowbreak{}0.5 & \{"alpha": 0.01, "l1": 0.5\} & 1.0446 \\
en\_\allowbreak{}0.1\_\allowbreak{}0.1 & \{"alpha": 0.1, "l1": 0.1\} & 1.0173 \\
en\_\allowbreak{}0.1\_\allowbreak{}0.5 & \{"alpha": 0.1, "l1": 0.5\} & 0.9984 \\
pls\_\allowbreak{}2 & \{"components": 2\} & 1.0056 \\
pls\_\allowbreak{}5 & \{"components": 5\} & 1.0271 \\
pls\_\allowbreak{}10 & \{"components": 10\} & 1.0460 \\
rf\_\allowbreak{}30 & \{"leaf": 30\} & 0.9979 \\
rf\_\allowbreak{}60 & \{"leaf": 60\} & 0.9969 \\
\end{longtable}

\endgroup

注。所有配置使用相同的全开发集五折。分数为在各折和各结局间等权平均的宏平均归一化 RMSE，越低越好。RF 均有 256 棵树；原共享 RF 保留最小叶节点样本量 15 和全部特征。新增独立 RF 每次划分使用一半特征。SVR 使用 RBF 核，gamma = scale，epsilon = 0.1。HGB 使用 7 个叶节点、100 次迭代、0.05 的学习率且不早停。弹性网最多迭代 10,000 次，容差为 0.000001。PLS 使用折内拟合的内部缩放。这些分数用于选择最终候选，不是对整个搜索流程性能的无偏估计。

\Needspace{8\baselineskip}

\par\medskip\noindent{\bfseries 表 S5 探索性选定模型在此前已查看留出集上的表现}\par\nopagebreak

\begingroup\singlespacing\fontsize{10}{12}\selectfont\setlength{\parindent}{0pt}

\begin{longtable}[]{@{}
  >{\raggedright\arraybackslash}p{(\linewidth - 6\tabcolsep) * \real{0.2200}}
  >{\raggedright\arraybackslash}p{(\linewidth - 6\tabcolsep) * \real{0.2700}}
  >{\raggedright\arraybackslash}p{(\linewidth - 6\tabcolsep) * \real{0.3400}}
  >{\raggedright\arraybackslash}p{(\linewidth - 6\tabcolsep) * \real{0.1700}}@{}}
\toprule\noalign{}
\begin{minipage}[b]{\linewidth}\raggedright
结局
\end{minipage} & \begin{minipage}[b]{\linewidth}\raggedright
Pearson r［95\% CI］
\end{minipage} & \begin{minipage}[b]{\linewidth}\raggedright
RMSE［95\% CI］
\end{minipage} & \begin{minipage}[b]{\linewidth}\raggedright
决定系数
\end{minipage} \\
\midrule\noalign{}
\endhead
\bottomrule\noalign{}
\endlastfoot
消极情绪 & 0.245 {[}0.055, 0.417{]} & 0.1661 {[}0.1461, 0.1871{]} & 0.0200 \\
积极情绪 & 0.332 {[}0.149, 0.493{]} & 0.1645 {[}0.1415, 0.1886{]} & 0.0447 \\
生活满意度 & 0.164 {[}\mbox{-0.029}, 0.346{]} & 0.2023 {[}0.1821, 0.2223{]} & 0.0155 \\
\end{longtable}

\endgroup

注。n = 104；独立随机森林，叶节点样本量 60、256 棵树，每次划分使用一半特征。Pearson 区间采用 Fisher z，RMSE 区间使用 5,000 次参与者自助抽样。单位仍为 CSV 保存单位。此前对结果的解释已经参考过该留出集，因此这些是探索性的条件区间，不是对修订后研究过程的独立验证。

\Needspace{8\baselineskip}

\par\medskip\noindent{\bfseries 表 S6 探索性配对误差比较}\par\nopagebreak

\begingroup\singlespacing\fontsize{10}{12}\selectfont\setlength{\parindent}{0pt}

\begin{longtable}[]{@{}
  >{\raggedright\arraybackslash}p{(\linewidth - 8\tabcolsep) * \real{0.2400}}
  >{\raggedright\arraybackslash}p{(\linewidth - 8\tabcolsep) * \real{0.2100}}
  >{\raggedright\arraybackslash}p{(\linewidth - 8\tabcolsep) * \real{0.1500}}
  >{\raggedright\arraybackslash}p{(\linewidth - 8\tabcolsep) * \real{0.2600}}
  >{\raggedright\arraybackslash}p{(\linewidth - 8\tabcolsep) * \real{0.1400}}@{}}
\toprule\noalign{}
\begin{minipage}[b]{\linewidth}\raggedright
比较
\end{minipage} & \begin{minipage}[b]{\linewidth}\raggedright
结局
\end{minipage} & \begin{minipage}[b]{\linewidth}\raggedright
RMSE 差值
\end{minipage} & \begin{minipage}[b]{\linewidth}\raggedright
95\% CI
\end{minipage} & \begin{minipage}[b]{\linewidth}\raggedright
p（Holm）
\end{minipage} \\
\midrule\noalign{}
\endhead
\bottomrule\noalign{}
\endlastfoot
新模型 vs 均值 & 消极情绪 & \mbox{-0.00341} & {[}\mbox{-0.00670}, \mbox{-0.00038}{]} & 0.168 \\
新模型 vs 均值 & 积极情绪 & \mbox{-0.00516} & {[}\mbox{-0.00819}, \mbox{-0.00215}{]} & 0.011 \\
新模型 vs 均值 & 生活满意度 & \mbox{-0.00252} & {[}\mbox{-0.00643}, 0.00137{]} & 0.817 \\
新模型 vs 原 RF & 消极情绪 & \mbox{-0.00002} & {[}\mbox{-0.00225}, 0.00223{]} & 1.000 \\
新模型 vs 原 RF & 积极情绪 & \mbox{-0.00040} & {[}\mbox{-0.00254}, 0.00161{]} & 1.000 \\
新模型 vs 原 RF & 生活满意度 & \mbox{-0.00009} & {[}\mbox{-0.00291}, 0.00266{]} & 1.000 \\
\end{longtable}

\endgroup

注。n = 104。差值为新选定模型减去相应参照，负值有利于新模型。区间使用 5,000 次配对参与者自助抽样；P 值使用 10,000 次预测交换，并对本扩展中的六项探索性比较作 Holm 校正。与原森林之差的所有区间均包含零。积极情绪与均值基线的比较达到本扩展内部的校正标准，但这不能消除探索性质，也不能改变嵌套开发集未见改善的结果。原始文件保留复用比较函数的标签；model\_\allowbreak{}a 和 model\_\allowbreak{}b 指明实际流程。另行从原始数据重新拟合后，全部修订模型预测完全复现（最大绝对差为 0）。

\clearpage

\section{补充图}\label{ux8865ux5145ux56fe}

\begin{figure}[!htbp]\centering
\includegraphics[width=\linewidth,height=0.54\textheight,keepaspectratio]{figures/figureS1_target_distributions.pdf}
\par\vspace{5pt}{\raggedright\singlespacing\fontsize{10}{12}\selectfont\noindent 图 S1. 1,037 名参与者最早保留记录中的主分析分数分布。使用 CSV 中存储的分数单位，未假定任何逆变换。组距为 .05，计数单位为参与者人数。\par}
\end{figure}

\begin{figure}[!htbp]\centering
\includegraphics[width=\linewidth,height=0.54\textheight,keepaspectratio]{figures/figureS2_sensitivity.pdf}
\par\vspace{5pt}{\raggedright\singlespacing\fontsize{10}{12}\selectfont\noindent 图 S2. 原开发阶段选定模型在主分析数据集及三项预定敏感性数据集上的表现。点表示 RMSE，误差线表示 CSV 单位下的条件性 95\% 参与者 bootstrap 区间。留出样本量见表 S2。模型使用固定的主分析参数重新拟合；该图不构成新的模型选择分析。\par}
\end{figure}

\end{document}
